\scanpage{47}

Notons que si $H$ est un gpe abstrait, $\mathrm{Tr}(H)$ l'ens des $(\rho,\rho',\rho'')\in H^3$ tels que $\rho\rho'\rho''=1$, le groupe $\unsure{\mathfrak S_3|\mathcal S_3}$ opère dedans de façon évidente, l'orbite de $(\rho,\rho',\rho'')$ \deleted{$(\rho,\rho',\rho'')\ \unsure{\text{et}}\ (\rho',\rho'',\unsure{\rho''|\rho})$} étant formée de

\[
(\rho,\rho',\rho'')\quad(\rho',\rho'',\rho)\quad(\rho'',\rho,\rho')\quad
(\rho''{}^{-1},\rho'{}^{-1},\rho^{-1})\quad
(\rho'{}^{-1},\rho^{-1},\rho''{}^{-1})\quad
(\rho^{-1},\rho''{}^{-1},\rho'{}^{-1})
\]

De cette façon, le gpe $\mathfrak S_3$ opère sur la cat. des ``groupes pythagoriciens avec présentation de Schwartz'' ou simplement Systèmes de Schwartz. On est intéressé à tenir compte de cette opération, et notamment déterminer quels sont les transformés \inserted{(par op. de $\mathfrak S_3$)} d'un $(H,\rho,\rho',\rho'')$ qui lui sont isom. \deleted{\unsure{et}}

\uline{Rmq} La notion d'isom qui nous intéresse est celle définie par le sch. en gpes ambiant $A$, ce qui s'explicite en \unsure{terme|termes} élémentaire ainsi :

a) Cas $\unsure{A_4|\mathcal A_4|\mathfrak A_4}$ : isomorphie par op. de $\mathfrak S_4$ \inserted{sur $A_4$} (\uuline{pas} $A_4$) On peut néanmoins se proposer de faire aussi la classification \uline{plus fine} \deleted{\unsure{du|dans}} d'op. de $A_4$.

b) Cas $\mathfrak S_4$ : isom par autom. int.

c) Cas $A_5$ : isom par autom. int. Il y a une classification moins fine, modulo autom. de $\mathfrak S_5$ sur $A_5$, qui donne ``moitié moins'' de classes d'isomorphie.

% archive sheet number: 46
